\documentclass[a4paper,11pt]{article}
\usepackage{lmodern}
\usepackage[T1]{fontenc}
\usepackage[utf8]{inputenc}
\usepackage{amsmath,amssymb,amsthm}
\usepackage{mathtools}
\usepackage{booktabs}
\usepackage{longtable}
\usepackage{array}
\usepackage{hyperref}
\usepackage{graphicx}
\usepackage{float}
\usepackage{geometry}
\geometry{margin=25mm}

\hypersetup{
  colorlinks=true,
  linkcolor=blue,
  urlcolor=blue,
  citecolor=blue
}

\theoremstyle{definition}
\newtheorem{theorem}{Theorem}[section]
\newtheorem{proposition}[theorem]{Proposition}
\newtheorem{lemma}[theorem]{Lemma}
\newtheorem{corollary}[theorem]{Corollary}
\newtheorem{definition}[theorem]{Definition}
\newtheorem{remark}[theorem]{Remark}
\newtheorem{observation}[theorem]{Observation}

\setlength{\parindent}{1em}
\setlength{\parskip}{0.5em}

\providecommand{\tightlist}{\setlength{\itemsep}{0pt}\setlength{\parskip}{0pt}}
\providecommand{\real}[1]{#1}
\begin{document}

\section{A Quantitative Correspondence between the Lattice Volume
Deficit and the Tangherlini Bekenstein--Hawking Entropy (Paper 4:
Entropy
Correspondence)}\label{a-quantitative-correspondence-between-the-lattice-volume-deficit-and-the-tangherlini-bekensteinhawking-entropy-paper-4-entropy-correspondence}

\textbf{Author}: Noriaki Kihara (WF System Co., Ltd.) \textbf{ORCID}:
\href{https://orcid.org/0009-0004-6753-4020}{0009-0004-6753-4020}
\textbf{Date}: April 2026 \textbf{DOI}:
\href{https://doi.org/10.5281/zenodo.19837594}{10.5281/zenodo.19837594}

\begin{center}\rule{0.5\linewidth}{0.5pt}\end{center}

\subsection{Abstract}\label{abstract}

We propose to interpret the volume deficit \(\Delta(R) = V_4(R) - N(k)\)
--- derived in a companion paper as the difference between the volume of
a four-dimensional ball and the count of integer-centred unit cubes
contained in it --- as a count of degrees of freedom in Planck units,
and to compare this count with the Bekenstein--Hawking entropy of a
four-plus-one-dimensional Schwarzschild--Tangherlini black hole. The two
quantities exhibit identical \(R^3\) scaling in \(R\), with the
asymptotic ratio \(\Delta(R) / S_{BH} \to 32/(3\pi) \approx 3.40\). The
correspondence is structural and quantitative; we leave the physical
interpretation of the constant ratio as an open problem.

\begin{center}\rule{0.5\linewidth}{0.5pt}\end{center}

\subsection{§1. Introduction}\label{introduction}

The Bekenstein--Hawking entropy of a black hole --- proportional to the
horizon area in Planck units --- admits an obvious dimensional analogue
in higher dimensions. For a four-plus-one-dimensional
Schwarzschild--Tangherlini black hole, the horizon is a three-sphere of
radius \(r_h\), and the Bekenstein--Hawking entropy is
\[S_{BH} = \frac{A_3}{4 G_5} = \frac{2\pi^2 r_h^3}{4 G_5} = \frac{\pi^2 r_h^3}{2 G_5}.\]

The entropy scales as \(r_h^3\). The microscopic origin of this entropy
--- what counts as a microstate --- remains an open problem in the
standard formulation.

In a separate strand of work, we have considered the geometric problem
of packing unit cubes inside a four-dimensional ball \(B(R)\) of radius
\(R\). The number \(N(k)\) of integer-centred unit cubes contained in
\(B(R)\) for \(R = 2k+1\) can be computed exactly; the volume deficit
\[\Delta(R) := V_4(R) - N(k), \qquad V_4(R) = (\pi^2/2) R^4,\] admits
the asymptotic expansion
\[\Delta(R) = \frac{16\pi}{3} R^3 - 6\pi R^2 + O(R), \qquad c := \lim_{k\to\infty} \frac{\Delta(R)}{2\pi^2 R^3} = \frac{8}{3\pi}.\]

The quantity \(\Delta(R)\) is a purely combinatorial quantity defined by
the discrete packing problem. In this paper we identify \(\Delta(R)\),
expressed in Planck units, with a candidate count of ``discrete drift
degrees of freedom'' associated with a four-plus-one-dimensional black
hole, and compare with \(S_{BH}\).

The principal observations are:

\begin{enumerate}
\def\labelenumi{\arabic{enumi}.}
\tightlist
\item
  Both \(\Delta(R)\) and \(S_{BH}\) scale as the cube of the radius.
\item
  Under the identification \(R \leftrightarrow r_h\), the ratio
  \(\Delta(R)/S_{BH}\) approaches the constant
  \(32/(3\pi) \approx 3.40\).
\item
  This ratio is independent of any free parameter in the construction;
  it is fixed by geometry and by the elementary identification of \(R\)
  with \(r_h\).
\end{enumerate}

The status of these observations is structural. We do not claim that
\(\Delta(R)\) \emph{equals} the Bekenstein--Hawking entropy. We claim
that \(\Delta(R)\) has the right scaling and a definite proportionality
to \(S_{BH}\), which is a non-trivial check on the discrete drift
hypothesis introduced elsewhere in the program.

The paper is organized as follows. Section 2 reviews the necessary
inputs from the companion papers. Section 3 sets up the identification
of \(R\) with the horizon radius \(r_h\) and computes the ratio
\(\Delta(R)/S_{BH}\). Section 4 discusses the physical interpretation of
the ratio. Section 5 concludes.

\begin{center}\rule{0.5\linewidth}{0.5pt}\end{center}

\subsection{§2. Inputs from Other
Papers}\label{inputs-from-other-papers}

\subsubsection{\texorpdfstring{§2.1 The volume deficit
\(\Delta(R)\)}{§2.1 The volume deficit \textbackslash Delta(R)}}\label{the-volume-deficit-deltar}

From the companion paper on the discrete packing problem, we have:
\[\Delta(R) = V_4(R) - N(k) = \frac{16\pi}{3} R^3 - 6\pi R^2 + 4R - 1 + O(R^{-1}) - E(R),\]
where \(E(R) = O(R^3)\) is the lattice-point error term, conjectured to
have zero leading-\(R^3\) coefficient on average.

The normalized leading constant is
\(c := \lim \Delta(R)/(2\pi^2 R^3) = 8/(3\pi) \approx 0.84883\).

\subsubsection{§2.2 The four-plus-one-dimensional Tangherlini
metric}\label{the-four-plus-one-dimensional-tangherlini-metric}

The Schwarzschild--Tangherlini metric in 4+1 dimensions is
\[ds^2 = -f_T(r)\, dt^2 + f_T(r)^{-1} dr^2 + r^2\, d\Omega_3^2, \qquad f_T(r) = 1 - \frac{8 G_5 M}{3\pi r^2}.\]
The horizon is at \(r_h = (8 G_5 M/(3\pi))^{1/2}\).

The Hawking temperature is \(T_H = 1/(2\pi r_h)\). The
Bekenstein--Hawking entropy is
\[S_{BH} = \frac{A_3}{4 G_5} = \frac{2\pi^2 r_h^3}{4 G_5}.\]

In Planck units defined by \(G_5 = \ell_P^3\) (a natural choice in 5
dimensions), this is \[S_{BH} = \frac{\pi^2 r_h^3}{2 \ell_P^3}.\]

\subsubsection{§2.3 Identification}\label{identification}

The construction of the central-projection metric (Paper 2) introduces a
curvature radius \(R\) as the natural geometric scale of the
four-dimensional subjective chart. The dynamical analysis of Paper 5
establishes that this \(R\) scales as \(M^{1/2}\) for a black hole of
mass \(M\), in agreement with the Tangherlini horizon-radius scaling.

We therefore identify \(R \leftrightarrow r_h\) throughout this paper.
The ratio \(\Delta(R)/S_{BH}\) is then a dimensionless number depending
only on the geometry of the construction.

\begin{center}\rule{0.5\linewidth}{0.5pt}\end{center}

\subsection{§3. The Quantitative
Correspondence}\label{the-quantitative-correspondence}

\subsubsection{\texorpdfstring{§3.1 Both quantities scale as
\(R^3\)}{§3.1 Both quantities scale as R\^{}3}}\label{both-quantities-scale-as-r3}

From §2.1: \(\Delta(R) \sim (16\pi/3) R^3\). From §2.2 (with
\(r_h \leftrightarrow R\) and \(G_5 = \ell_P^3\)):
\(S_{BH} = \pi^2 R^3 / (2 \ell_P^3)\).

In Planck units (\(\ell_P = 1\)):
\[\Delta(R) \sim \frac{16\pi}{3} R^3, \qquad S_{BH} \sim \frac{\pi^2}{2} R^3.\]

Both quantities scale as the cube of the radius.

\subsubsection{§3.2 The asymptotic ratio}\label{the-asymptotic-ratio}

The ratio is
\[\frac{\Delta(R)}{S_{BH}} \to \frac{16\pi/3}{\pi^2/2} = \frac{32}{3\pi} \approx 3.397.\]

This is a dimensionless constant, independent of \(R\).

\subsubsection{\texorpdfstring{§3.3 Numerical values at finite
\(k\)}{§3.3 Numerical values at finite k}}\label{numerical-values-at-finite-k}

The convergence to the asymptotic ratio is observed in the data:

{\def\LTcaptype{none} % do not increment counter
\begin{longtable}[]{@{}rrrrr@{}}
\toprule\noalign{}
\(k\) & \(R\) & \(\Delta(R)\) & \(S_{BH} = (\pi^2/2) R^3\) &
\(\Delta/S_{BH}\) \\
\midrule\noalign{}
\endhead
\bottomrule\noalign{}
\endlastfoot
10 & 21 & 141,580.27 & 45,712.2 & 3.10 \\
20 & 41 & 1,114,426.60 & 340,184.3 & 3.28 \\
30 & 61 & 3,708,117.64 & 1,119,679.5 & 3.31 \\
40 & 81 & 9,064,209.71 & 2,624,489.6 & 3.45 \\
50 & 101 & 16,981,622.84 & 5,054,121.0 & 3.36 \\
60 & 121 & 29,303,164.67 & 8,740,869.3 & 3.35 \\
\end{longtable}
}

The ratio approaches \(32/(3\pi) \approx 3.40\) from below, with
finite-\(k\) deviations consistent with the sub-leading \(R^2\)
correction to \(\Delta(R)\) (which contributes to the lower numerical
values at small \(k\)).

\subsubsection{§3.4 Status of the
correspondence}\label{status-of-the-correspondence}

The correspondence in scaling and the constant value of the ratio is
established. The physical interpretation of the ratio \(32/(3\pi)\) is
open. Possibilities include:

\begin{enumerate}
\def\labelenumi{\arabic{enumi}.}
\item
  \emph{Geometric.} The ratio \(32/(3\pi)\) relates the boundary
  correction \(16\pi/3\) (from inclusion--exclusion on the auxiliary
  body \(\Omega_R\)) to the horizon area density \(\pi^2/2\) (from
  \(V_4(R) = (\pi^2/2) R^4\) and \(S_3(R) = 2\pi^2 R^3\)). Specifically:
  \[\frac{32}{3\pi} = \frac{16\pi/3}{\pi^2/2} = \frac{2 \cdot (16\pi/3)}{2\pi^2} = \frac{2 c \cdot 2\pi^2}{2\pi^2} = 2 c \cdot \frac{2\pi^2}{\pi^2} = 4c \approx 4 \cdot 0.849 \approx 3.40.\]
  Equivalently, the ratio is \(4c\) where \(c = 8/(3\pi)\) is the
  volume-deficit constant.
\item
  \emph{Counting normalization.} The factor of \(4\) in
  \(S_{BH} = A/(4 \ell_P^{D-2})\) is the Bekenstein normalization (one
  quarter of the area in Planck units). A different counting convention
  would shift the ratio by a factor.
\item
  \emph{Dimensional convention.} The identification \(G_5 = \ell_P^3\)
  is one natural choice; other choices (\(G_5 = \alpha \ell_P^3\) for
  some \(\alpha\)) shift \(S_{BH}\) and hence the ratio.
\end{enumerate}

We do not at this stage privilege any of these interpretations.

\begin{center}\rule{0.5\linewidth}{0.5pt}\end{center}

\subsection{§4. Physical Interpretation: Open
Questions}\label{physical-interpretation-open-questions}

\subsubsection{\texorpdfstring{§4.1 Is \(\Delta(R)\) ``the''
entropy?}{§4.1 Is \textbackslash Delta(R) ``the'' entropy?}}\label{is-deltar-the-entropy}

We have shown that \(\Delta(R)\) has the same scaling as \(S_{BH}\) and
is in a constant ratio to it. This does not establish that
\(\Delta(R) = S_{BH}\) in any deeper sense. The two are quantities of
the same dimensional type (counts of degrees of freedom in Planck units,
both scaling as length cubed in 4+1 dimensions), and they bear a
structural relation.

A stronger interpretation would require: (i) An independent derivation
of \(\Delta(R)\) from a microscopic theory of black hole microstates.
(ii) A consistent extension to logarithmic and finite-size corrections.
(iii) A clear physical mechanism for the discrete drift that justifies
the identification.

The companion papers (3 and 5) address (i) and (iii) at the level of
geometric and combinatorial structure, but not at the level of a
complete microscopic theory.

\subsubsection{\texorpdfstring{§4.2 The constant
\(32/(3\pi) \approx 3.40\)}{§4.2 The constant 32/(3\textbackslash pi) \textbackslash approx 3.40}}\label{the-constant-323pi-approx-3.40}

The ratio is a definite number. Its physical interpretation, if any, is
open. It might be a manifestation of:

\begin{itemize}
\tightlist
\item
  A residual gauge of the discrete-cube packing problem (the choice of
  integer lattice rather than another regular lattice).
\item
  A specific choice of dimensional reduction from a putative
  higher-dimensional structure.
\item
  An artifact of the inclusion--exclusion boundary correction having a
  specific geometric interpretation.
\end{itemize}

No preferred interpretation is offered here.

\subsubsection{§4.3 Comparison with logarithmic and one-loop
corrections}\label{comparison-with-logarithmic-and-one-loop-corrections}

The standard Bekenstein--Hawking entropy receives logarithmic and
one-loop corrections in the semi-classical approximation. The
volume-deficit construction has its own sub-leading corrections --- the
\(-6\pi R^2\) term and lower --- which we have not attempted to match
against the standard quantum corrections. \textbf{後で検討}: detailed
correspondence at sub-leading orders.

\begin{center}\rule{0.5\linewidth}{0.5pt}\end{center}

\subsection{§5. Conclusion}\label{conclusion}

We have identified a structural correspondence between the volume
deficit \(\Delta(R)\), derived from the geometric problem of discrete
unit-cube packing, and the Bekenstein--Hawking entropy of a
four-plus-one-dimensional Schwarzschild--Tangherlini black hole. Both
quantities scale as \(R^3\), with the asymptotic ratio
\(\Delta(R)/S_{BH} \to 32/(3\pi) \approx 3.40\) in Planck units.

This is a non-trivial check on the discrete drift hypothesis: under the
assumption that \(\Delta(R)\) is the relevant entropy quantity for the
discrete drift, the scaling matches the expected form. The
identification of the constant ratio as something physical, rather than
as a feature of the geometric construction, is left as an open problem.

The dynamical determination of \(R\) from the matter content of the
four-dimensional theory is the subject of Paper 5.

\begin{center}\rule{0.5\linewidth}{0.5pt}\end{center}

\subsection{References (selected,
後で検討)}\label{references-selected-ux5f8cux3067ux691cux8a0e}

\begin{itemize}
\tightlist
\item
  J. D. Bekenstein, ``Black Holes and Entropy,'' \emph{Phys. Rev.~D} 7,
  2333 (1973).
\item
  S. W. Hawking, ``Particle Creation by Black Holes,'' \emph{Comm. Math.
  Phys.} 43, 199 (1975).
\item
  F. R. Tangherlini, ``Schwarzschild Field in \(n\) Dimensions and the
  Dimensionality of Space Problem,'' \emph{Nuovo Cimento} 27, 636
  (1963).
\item
  R. C. Myers and M. J. Perry, ``Black Holes in Higher Dimensional
  Spacetimes,'' \emph{Annals Phys.} 172, 304 (1986).
\end{itemize}

\end{document}
